
\section{Dvigubos grupės}

Šiame skyriuje sujungsime pirmų dviejų skyrių idėjas.
Pasinaudosime dvigubu atvaizdavimu $\rho$,
kad perkelti taškinę grupę $G$ iš $\Orth3$ į $\pinp$ arba $\pinm$.
Naujai gauta grupė yra vadinama grupės $G$ dviguba grupe.
\begin{definition}
	Grupės $\Orth3$ pogrupio $G$ \emph{dviguba grupė viduje grupės $\pinp$} yra
	$G$ pirmavaizdis (ang. \emph{preimage}) $\rho^{-1}(G)$ ir yra žymima $2G^+$.
	\emph{Dviguba grupė viduje grupės $\pinm$} yra
	$G$ pirmavaizdis $\rho^{-1}(G)$ ir yra žymima $2G^-$.
	Kai $G \in \SO3$, grupės $2G^+$ ir $2G^-$ yra izomorfiškos,
	nes $\pinp$ ir $\pinm$ lyginiai pogrųpiai yra izomorfiški,
	ir dvigubą grupę žymime $2G$.
\end{definition}
\begin{remark}
	Dvigubas grupes taip pat galima apibėžti
	kaip grupes $2G^+$ ir $2G^-$ tokias kad diagramos
	\begin{equation}
		\begin{tikzcd}
			\pinp \arrow[r, "\rho"] & \Orth3 \\
			2G^+ \arrow[u, "\iota", hook] \arrow[r, "\rho"] & G \arrow[u, "\iota", hook]
		\end{tikzcd}
		\quad
		\begin{tikzcd}
			\pinm \arrow[r, "\rho"] & \Orth3 \\
			2G^- \arrow[u, "\iota", hook] \arrow[r, "\rho"] & G \arrow[u, "\iota", hook]
		\end{tikzcd}
	\end{equation}
	yra komutatyvios.
	Kategorijų teorijoje ši konstrukcija yra vadinama
	grupių $G$ ir $\pinpm$ atgaline sankaba (ang. \emph{pullback})
	per homomorfizmus $\rho$ ir $\iota$.
\end{remark}

Aptarkime kaip galima dvigubos grupės prezentaciją.
Tegul $G = \present{ a,b,c,\cdots }{ r'=r,\cdots }$.
Atvaizdavimo $\rho$ branduolys yra $\qty{\mathbb1,-\mathbb1}$,
tad $-\mathbb1$ priklauso $2G$,
prezentacijų kontekste šį elementą žymėsime $z:=-\mathbb1$.
Elementą $z$ įprasta laikyti \emph{pasukimu kampu $2\pi$},
nes į $\mathup{SU}(2)$ pasukimo aplink vektorių $n$ kampu $\theta$ formulę
$e^{\frac{i\theta}{2}(\vb n,\symbf\sigma)}$
įstačius $\theta = 2\pi$ gauname $-I$.
Generatoriui $a$ jo pirmavaizdžio 
$\rho^{-1}(a) = \qty{ v \in \pinpm \mid \rho(v)=a }$ 
pasirinkta elementą pažymime $\tilde a$,
kitas  elementas turi būti $z\tilde a$.
Tokiu pat būdu pažymime $\tilde b,\tilde c,\cdots$.
\begin{proposition}
	Elementai $z, \tilde a, \tilde b, \tilde c, \cdots \in \pinpm$ generuoja grupę $2G$.
\end{proposition}
\begin{proof}
	Elementas $g \in G$ yra žodis išrekštas raidėmis $a,b,c,\cdots$.
	Tegul $\tilde g$ yra žodis gaunamas žodyje $g$ raides $a, b, c \cdots$
	pakeičiant į $\tilde a, \tilde b, \tilde c, \cdots$.
	Turime kad $\rho(\tilde g) = g$ ir $\rho(z\tilde g)=g$,
	nes \rho yra homomorfizmas, $\rho(z)=e$ ir visiem gerneratoriams $\rho(\tilde a)=a$.
	Elementai $\tilde g$ ir $z\tilde g$ sudaro elemento $g$ primavaizdį
	ir $g$ gali būti bet koks grupės $G$ elementas,
	tad raidės $z, \tilde a, \tilde b, \tilde c, \cdots \in \pinpm$ generuoja visą $2G$.
\end{proof}
Su sąryšiasi yra sudetingesnė situacija.
Automatiškai turime sąryšį $z^2=e$ ir sąryšius $z\tilde a =\tilde a z$, $z\tilde b =\tilde b z$, $\cdots$,
o pergeliant grupės $G$ apibriažiančius sąryšius atsiranda dvi galimybės:
vienas sąryšis $r'=r$ grupėje $G$,
gali virsti į sąryšį $\tilde r'=\tilde r$
arba į sąryšį $\tilde r'=z\tilde r$ grupėje $2G$.
Elemetų $\tilde r'$ ir $\tilde r$ išraiškos yra fiksuotos pasirinkimo
kuriuos elementus žymime $\tilde a, \tilde b, \tilde c, \cdots$,
tad kuris iš dviejų sąryšiu yra teisingas reikia patikrinti sudauginus
generatorius grupėje $\pinpm$ ir supaprastinant išrišką su $\pinpm$ sąvybėmis.



\subsection{Chiralinių grupų dvigubos grupės}


\subsection{Achiralinių grupių paritetinės dvigubos grupės}


\subsection{Achiralinių grupių apariterinės grupės}


\clearpage{}
